\documentclass[a4paper,12pt]{article}
\usepackage[slovene]{babel}
\usepackage[utf8]{inputenc}
\usepackage{amsmath,amssymb}
\usepackage[margin=2cm]{geometry}
\usepackage{lmodern}
\usepackage{titlesec}
\usepackage{xcolor}
\usepackage{tikz}
\usepackage{hyperref}  
\usepackage{cleveref}
\usepackage{pgfplots}
\pgfplotsset{compat=1.18}
% ==========================================
% Avtomatsko štetje nalog po sekcijah
% ==========================================

\usepackage{etoolbox}

\newcounter{secitemcount}      % števec itemov v sekciji
\newcounter{totalitemcount}    % skupni števec vseh nalog

\newcommand{\allresults}{}     % shranjevanje rezultatov

% ob začetku vsake sekcije resetiramo lokalni števec
\pretocmd{\section}{%
  \setcounter{secitemcount}{0}%
  \gdef\currentsectiontitle{}%
}{}{}

% prestrežemo naslov sekcije
\pretocmd{\section}{\gdef\currentsectiontitle{#1}}{}{}

% vsak \item poveča oba števca
\pretocmd{\item}{%
  \stepcounter{secitemcount}%
  \stepcounter{totalitemcount}%
}{}{}

% ob koncu enumerate avtomatsko shranimo rezultat
\AtEndEnvironment{enumerate}{%
  \gappto{\allresults}{%
    \resultentry{\currentsectiontitle}{\arabic{secitemcount}}%
  }%
}

% kako se izpiše posamezen rezultat v Rešitvah
\newcommand{\resultentry}[2]{%
  \subsection*{#1 (število nalog: #2)}%
}

%%%%%%%%%%%%%%%%%%%%%%%%%%%%%%%%%%%%%%%%%%%%%%%%%%%
% Definicija rjave barve
\definecolor{mybrown}{RGB}{140, 90, 40}

% SECTION – rjav okvir, bela notranjost
\titleformat{\section}
  {\large\bfseries}
  {%
    \fcolorbox{mybrown}{white}{\rule{1.2ex}{1.2ex}}%
    \quad\thesection
  }
  {0.8em}
  {}

% SUBSECTION – manjši rjav okvir, bela notranjost
\titleformat{\subsection}
  {\bfseries}
  {%
    \fcolorbox{mybrown}{white}{\rule{0.9ex}{0.9ex}}%
    \quad\thesubsection
  }
  {0.8em}
  {}

% --- Naslov dokumenta z rjavim kvadratkom ---
\newcommand{\DocTitle}[1]{%
  \begin{center}
    {\rule{\textwidth}{0.1mm}\\[0.4cm]
    \LARGE
    \fcolorbox{mybrown}{gray!40}{\rule{1.4ex}{1.4ex}}
    \quad #1
    }
  \end{center}
  \vspace{1em}
}




\begin{document}

\pagestyle{fancy}
\lhead{}\chead{}\rhead{}

\begin{center}
\Huge {\textbf{Test – Odvod funkcije}}
\end{center}


% ===================== 1 =====================
\subsection{ Izračunaj odvod funkcije:}
\def\tockeA{2+2}
\addtocounter{tocke}{\tockeA}


a) $f(x)=5x^3-2x^2+7x-3$

b) $g(x)=\dfrac{5}{x^3}$

\vspace{1cm}

% ===================== 2 =====================


\addtocounter{tocke}{2+3}
\subsection{ Izračunaj odvod v dani točki:}

a) $f(x)=x^3-3x^2+2x$, \quad $x_0=3$

b) $g(x)=\frac{3x+2}{x-2}$, \quad $x_0=4$
\end{naloga}

\vspace{1cm}

% ===================== 3 =====================
\def\tockeC{2+2+2}
\addtocounter{tocke}{2+2+2}
\subsection{ Določi enačbo tangente na graf funkcije v dani točki:}

a) $f(x)=x^2$, \quad $x_0=2$

b) $g(x)=\sqrt{x}$, \quad $x_0=9$

c) $h(x)=\frac{1}{x}$, \quad $x_0=2$

\vspace{1cm}

% ===================== 4 =====================

\addtocounter{tocke}{3+3}
\subsection{Določi kot, ki ga tangenta na graf funkcije v dani točki oklepa z osjo $x$:}

a) $f(x)=\frac{5}{x}$, \quad $x_0=2$

b) $g(x)=\dfrac{1}{4}x^3$, \quad $x_0=3$


\vspace{1cm}

% ===================== 5 =====================

\addtocounter{tocke}{3+3}
\subsection{Izračunaj kot med krivuljama v presečišču:}

a) $f(x)=x^2$, \quad $g(x)=x^2-4x$

b) $f(x)=x^2+2$, \quad $g(x)=-x^2+6$

\vfill


\newpage

% ===================== REŠITVE =====================
\section*{Rešitve}

\subsection*{1}
a)
\[
f'(x)=15x^2-4x+7
\]

b)
\[
g(x)=5x^{-3} \Rightarrow g'(x)=-15x^{-4}=-\frac{15}{x^4}
\]

\subsection*{2}
a)
\[
f'(x)=3x^2-6x+2
\]
\[
f'(3)=27-18+2=11
\]

b)
\[
g'(x)=\frac{3(x-2)-(3x+2)}{(x-2)^2}
=\frac{3x-6-3x-2}{(x-2)^2}
=\frac{-8}{(x-2)^2}
\]
\[
g'(4)=\frac{-8}{4}=-2
\]

\subsection*{3}
Uporabimo: $y=f'(x_0)(x-x_0)+f(x_0)$

a)
\[
f'(x)=2x,\quad f'(2)=4,\quad f(2)=4
\]
\[
y=4(x-2)+4=4x-4
\]

b)
\[
g'(x)=\frac{1}{2\sqrt{x}},\quad g'(9)=\frac{1}{6},\quad g(9)=3
\]
\[
y=\frac{1}{6}(x-9)+3
\]

c)
\[
h'(x)=-\frac{1}{x^2},\quad h'(2)=-\frac{1}{4},\quad h(2)=\frac{1}{2}
\]
\[
y=-\frac{1}{4}(x-2)+\frac{1}{2}=-\frac{1}{4}x+1
\]

\subsection*{4}
Velja: $\tan \alpha = f'(x_0)$

a)
\[
f'(x)=-\frac{5}{x^2},\quad f'(2)=-\frac{5}{4}
\]
\[
\alpha=\arctan\!\left(-\tfrac{5}{4}\right)
\]

b)
\[
g'(x)=\frac{3}{4}x^2,\quad g'(3)=\frac{27}{4}
\]
\[
\alpha=\arctan\!\left(\tfrac{27}{4}\right)
\]

\subsection*{5}
Formula:
\[
\tan \varphi = \left|\frac{k_1-k_2}{1+k_1k_2}\right|
\]

a)
Presečišče:
\[
x^2=x^2-4x \Rightarrow x=0
\]
\[
f'(x)=2x,\quad g'(x)=2x-4
\]
\[
k_1=0,\quad k_2=-4
\]
\[
\tan \varphi = 4
\]

b)
Presečišče:
\[
x^2+2=-x^2+6 \Rightarrow 2x^2=4 \Rightarrow x=\pm\sqrt{2}
\]

Za $x=\sqrt{2}$:
\[
f'(x)=2x=2\sqrt{2},\quad g'(x)=-2x=-2\sqrt{2}
\]
\[
\tan \varphi = \left|\frac{2\sqrt{2}-(-2\sqrt{2})}{1-8}\right|=\frac{4\sqrt{2}}{7}
\]

(enak rezultat dobimo tudi za $x=-\sqrt{2}$)

\end{document}
